\documentclass[11pt]{article}

\usepackage[a4paper,margin=2.2cm]{geometry}
\usepackage[T1]{fontenc}
\usepackage[utf8]{inputenc}
\usepackage{lmodern}
\usepackage{amsmath,amssymb,mathtools}
\usepackage{microtype}
\usepackage[dvipsnames]{xcolor}
\usepackage[most]{tcolorbox}
\usepackage{enumitem}
\usepackage{titlesec}
\usepackage{booktabs}

% ---------- Section formatting ----------
\titleformat{\section}{\Large\bfseries}{\thesection}{0.75em}{}
\titleformat{\subsection}{\large\bfseries}{\thesubsection}{0.75em}{}
\titlespacing*{\section}{0pt}{1.4ex plus .2ex minus .2ex}{0.9ex}
\titlespacing*{\subsection}{0pt}{1.0ex plus .2ex minus .2ex}{0.5ex}

% ---------- Box environments ----------
\newtcolorbox{notebox}[1]{
    breakable, enhanced,
    colback=gray!4, colframe=gray!55,
    boxrule=0.7pt, arc=1.5mm,
    left=2.5mm, right=2.5mm, top=1.3mm, bottom=1.3mm,
    title={\bfseries #1}, coltitle=black,
    colbacktitle=gray!12, fonttitle=\bfseries
}

\newtcolorbox{resultbox}[1]{
    breakable, enhanced,
    colback=blue!2, colframe=blue!40!black,
    boxrule=0.7pt, arc=1.5mm,
    left=2.5mm, right=2.5mm, top=1.3mm, bottom=1.3mm,
    title={\bfseries #1}, coltitle=black,
    colbacktitle=blue!10, fonttitle=\bfseries
}

\newtcolorbox{warningbox}[1]{
    breakable, enhanced,
    colback=orange!4, colframe=orange!60!black,
    boxrule=0.7pt, arc=1.5mm,
    left=2.5mm, right=2.5mm, top=1.3mm, bottom=1.3mm,
    title={\bfseries #1}, coltitle=black,
    colbacktitle=orange!14, fonttitle=\bfseries
}

% ---------- Helper commands ----------
\newcommand{\R}{\mathbb{R}}
\newcommand{\audititem}[2]{\par\smallskip\noindent\textbf{#1.} #2\par}
\newcommand{\statusline}[1]{\texttt{#1}}

\begin{document}

\begin{center}
    {\LARGE\bfseries Verification Note for the Dual-Gantry CT LPV-LFR Reformulation}\\[0.5em]
    {\large Consolidated audit document based on Drenth Chapter 2 and the current gantry derivation}\\[0.3em]
    {\normalsize Working document for internal verification, not the final supervisor-facing version}
\end{center}

\vspace{0.7em}

%------------------------------------------------------------------------------
\section{Purpose and Scope}
%------------------------------------------------------------------------------

This document consolidates the current verification work for the dual-gantry
continuous-time LPV-LFR reformulation into one \LaTeX{} file. Its goal is
\textbf{traceability}, not compactness. The intended use is:

\begin{itemize}[leftmargin=1.2em]
    \item check what Drenth's thesis actually states,
    \item separate that from what is only a generalized modeling lesson,
    \item identify where the dual-gantry derivation becomes our own algebra,
    \item keep the well-posedness discussion honest by separating Drenth's
    generic sufficient route from the sharper plant-specific route.
\end{itemize}

\begin{notebox}{Source Backbone}
This verification note is consolidated from the following working notes and
companion derivations already present in the repo:
\begin{itemize}[leftmargin=1.2em]
    \item \texttt{docs/drenth/ch2-sec21-source.md}
    \item \texttt{docs/drenth/ch2-sec211-source.md}
    \item \texttt{docs/drenth/ch2-sec22-source.md}
    \item \texttt{docs/drenth/ch2-generalized-recipe.md}
    \item \texttt{docs/drenth/ch2-dual-gantry-mapping.md}
    \item \texttt{LPV/LFR-derivation.tex}
    \item \texttt{LPV/M-invertibility.tex}
\end{itemize}
It therefore uses the notes as an audited backbone rather than re-inventing the
interpretation from scratch.
\end{notebox}

\begin{notebox}{Reading Rule}
For every major verification step below, the same fields are kept explicit:
\begin{itemize}[leftmargin=1.2em]
    \item Drenth reference
    \item What Drenth does
    \item Generalized step
    \item Dual-gantry application
    \item Status
    \item What still needs checking
\end{itemize}
The point is not elegance but the ability to verify each step carefully.
\end{notebox}

\begin{warningbox}{Important Source Boundary}
The thesis source for the generic LPV-LFR definition used here is Drenth's
\textbf{thesis}, Chapter~2, especially Sections~2.1--2.2
\cite{drenth2025thesis}. The IFAC paper \cite{drenth2025lpvlfr} remains useful
as supporting context, but its formal LPV-LFR definition is in discrete time.
So this document treats the thesis as the primary CT source.
\end{warningbox}

\begin{notebox}{Role of the Sources}
The source roles in this document are intentionally separated:
\begin{itemize}[leftmargin=1.2em]
    \item Section~2.1 of Drenth gives the generic CT LPV-LFR framework.
    \item Section~2.1.1 of Drenth motivates the preference for rational
    dependency over an unnecessarily overbounded affine embedding.
    \item Section~2.2 of Drenth gives a generic sufficient well-posedness route.
    \item \texttt{LPV/LFR-derivation.tex} provides the dual-gantry-specific
    realization algebra.
    \item \texttt{LPV/M-invertibility.tex} provides the sharper
    plant-specific well-posedness proof used for the current baseline.
\end{itemize}
\end{notebox}

%------------------------------------------------------------------------------
\section{Starting Point: Dual-Gantry CT Model}
%------------------------------------------------------------------------------

The dual-gantry first-principles model is taken from the current gantry
derivation already present in the repo. That derivation in turn starts from the
logical-coordinate mechanical model of Garcia-Herreros et al.~\cite{garcia2013model}. In logical
coordinates
\[
q = [X,\Theta,Y]^\top \in \R^3,
\]
the equations of motion are
\begin{equation}
    M(Y)\,\ddot{q}(t) + C\,\dot{q}(t) + K\,q(t) = f_\ell(t),
    \label{eqv:eom}
\end{equation}
with
\begin{equation}
    M(Y) =
    \begin{bmatrix}
        m_1 + m_2 + m_b + m_h &
        \dfrac{(m_1-m_2)L_b}{2} - m_h Y &
        0 \\
        \dfrac{(m_1-m_2)L_b}{2} - m_h Y &
        J_b + J_h + \dfrac{(m_1+m_2)L_b^2}{4} + m_h d^2 + m_h Y^2 &
        -m_h d \\
        0 & -m_h d & m_h
    \end{bmatrix}.
    \label{eqv:M}
\end{equation}

With the state
\[
x = [q^\top,\dot{q}^\top]^\top \in \R^6
\]
and generalized logical input
\[
u := f_\ell,
\]
the continuous-time state-space form is
\begin{equation}
    \dot{x}(t) = A_c(Y)\,x(t) + B_c(Y)\,u(t),
    \qquad
    y(t) = C_c\,x(t),
    \label{eqv:ctss}
\end{equation}
where
\begin{equation}
    A_c(Y) =
    \begin{bmatrix}
        0 & I_3 \\
        -M(Y)^{-1}K & -M(Y)^{-1}C
    \end{bmatrix},
    \qquad
    B_c(Y) =
    \begin{bmatrix}
        0 \\
        M(Y)^{-1}
    \end{bmatrix},
    \qquad
    C_c = \begin{bmatrix} I_3 & 0 \end{bmatrix}.
    \label{eqv:ctmats}
\end{equation}

The key structural point is that the dependence on the scheduling coordinate
$Y$ enters through $M(Y)^{-1}$, so the exact continuous-time model is
\textbf{rational} in $Y$.

\begin{resultbox}{Why an LPV-LFR reformulation is needed}
The exact dual-gantry baseline is not naturally affine in the scheduling
variable $Y$. The LPV-LFR reformulation is therefore not a cosmetic rewriting:
it is the device that allows the rational dependence induced by
$M(Y)^{-1}$ to be represented inside a constant interconnection plus a
scheduling block.
\end{resultbox}

%------------------------------------------------------------------------------
\section{Drenth-Based Stepwise Verification}
%------------------------------------------------------------------------------

This section is the core audit trail. Its internal flow is:
\begin{itemize}[leftmargin=1.2em]
    \item Steps~1--3: fix the target LPV-LFR framework and the CT model that
    must be recovered after collapse.
    \item Steps~4--6: explain why a rational repeated-$\Delta$ structure is the
    natural modeling direction.
    \item Steps~7--10: carry out the dual-gantry-specific realization.
    \item Step~11: classify the collapsed result.
\end{itemize}
The well-posedness discussion is then handled separately in the next section so
that the two proof routes do not get mixed together.

\subsection{Step 1: Fix the target representation class as a CT LPV-LFR}

\audititem{Drenth reference}{Thesis Section~2.1, especially eq.~(2.1)
\cite{drenth2025thesis}.}

\audititem{What Drenth does}{Drenth defines an LPV-LFR in continuous time as a
pair $(G,\Delta(p))$ with the interconnection
\[
\dot{x} = A x + B_w w + B_u u,\quad
z = C_z x + D_{zw} w + D_{zu} u,\quad
y = C_y x + D_{yw} w + D_{yu} u,\quad
w = \Delta(p)\,z.
\]
(Drenth writes $A_x$ for the state matrix; we use $A$ since no ambiguity arises.)
The matrices in $G$ are constant, while scheduling dependence enters through
$\Delta(p)$.}

\audititem{Generalized step}{Before any plant-specific algebra, choose the CT
LPV-LFR pair $(G,\Delta(p))$ as the target representation class.}

\audititem{Dual-gantry application}{For the dual gantry, the target is a CT
LPV-LFR with state $x = [q^\top,\dot{q}^\top]^\top$, generalized logical input
$u = f_\ell$, and measured output $y = q$.}

\audititem{Status}{\statusline{Direct from Drenth} for the framework choice.
\statusline{Own derivation} for the particular gantry state, input, and output
assignment.}

\audititem{What still needs checking}{This step does not justify the existence
of a convenient dual-gantry realization. It only fixes the target class.}

\subsection{Step 2: Treat the model as a constant interconnection plus a scheduling block}

\audititem{Drenth reference}{Thesis Section~2.1, conceptual text around
eq.~(2.1) and the definition of $(G,\Delta(p))$ \cite{drenth2025thesis}.}

\audititem{What Drenth does}{Drenth presents the LPV-LFR as the interconnection
between a nominal constant system $G$ and a structured scheduling block
$\Delta(p)$.}

\audititem{Generalized step}{Organize the reformulation so that all constant
linear structure lives in $G$ and all scheduling dependence is pushed into
$\Delta(p)$.}

\audititem{Dual-gantry application}{In the dual-gantry derivation, the constant
objects are $M_0$, $M_1$, $M_2$, $C$, $K$, and the eventual blocks of $G$.
The scheduling object is the scalar coordinate $Y$. The design target is that
$Y$ enters only through the eventual $\Delta(Y)$.}

\audititem{Status}{\statusline{Direct from Drenth} for the structural split.
\statusline{Generalized from Drenth} for using this split as a design rule
during reformulation.}

\audititem{What still needs checking}{This step does not yet tell us how the
$Y$-dependence is actually pushed into the loop.}

\subsection{Step 3: State the collapsed CT model that the LPV-LFR must reproduce}

\audititem{Drenth reference}{Thesis Section~2.1, eqs.~(2.3)--(2.4)
\cite{drenth2025thesis}.}

\audititem{What Drenth does}{Drenth states that, after eliminating the latent
variables, the LPV-LFR becomes a continuous-time LPV state-space model with
rational dependence on the scheduling variable.}

\audititem{Generalized step}{Write down the continuous-time target model that
must be recovered after elimination of the latent loop.}

\audititem{Dual-gantry application}{The target is exactly the MATLAB-derived CT
model in \eqref{eqv:ctss}--\eqref{eqv:ctmats}. The LPV-LFR is only acceptable if
it collapses back to these equations.}

\audititem{Status}{\statusline{Generalized from Drenth} for using the collapsed
form as the central correctness criterion. \statusline{Own derivation} for the
particular dual-gantry target equations.}

\audititem{What still needs checking}{This step does not tell us how to
construct the realization. It only states what success means.}

\subsection{Step 4: Decide whether the exact realization should be affine or rational}

\audititem{Drenth reference}{Thesis Section~2.1.1 for the motivating discussion;
Section~2.1, paragraph after eq.~(2.4), for the affine/rational classification
\cite{drenth2025thesis}.}

\audititem{What Drenth does}{Drenth identifies affine dependency as the special
case $D_{zw}=0$ and then argues in Section~2.1.1 that rational dependency can
reduce overbounding and reduce unnecessary scheduling-variable count.}

\audititem{Generalized step}{Before constructing the realization, decide whether
the exact plant can be represented adequately in the affine special case, or
whether the model should retain genuinely rational scheduling dependence.}

\audititem{Dual-gantry application}{Because the dual-gantry CT matrices depend
on $M(Y)^{-1}$, the exact collapsed dependence is rational in $Y$. So the
natural target is a rational LPV-LFR, not the affine special case.}

\audititem{Status}{\statusline{Generalized from Drenth} for the modeling
decision. \statusline{Own derivation} for concluding that this specific plant is
rational in $Y$.}

\audititem{What still needs checking}{This step does not yet construct the
rational realization or prove that the chosen one is minimal.}

\subsection{Step 5: Choose the scheduling variable with coupling loss in mind}

\audititem{Drenth reference}{Thesis Section~2.1.1, especially the discussion
around eqs.~(2.6)--(2.9) \cite{drenth2025thesis}.}

\audititem{What Drenth does}{Drenth compares an affine MSD embedding using two
scheduling variables $(p_1=x,\; p_2=x^2)$ with a rational embedding using only
one scheduling variable repeated in $\Delta$. The point is that the affine
embedding loses coupling between quantities that are not truly independent.}

\audititem{Generalized step}{Choose the scheduling description so that real
structural coupling is preserved as much as possible, instead of introducing
unnecessary independent scheduling variables.}

\audititem{Dual-gantry application}{For the gantry, the natural physical
scheduling quantity is the single coordinate $Y$. Even though $Y^2$ appears in
$M(Y)$, it is not introduced as a second independent scheduling variable.}

\audititem{Status}{\statusline{Generalized from Drenth}.}

\audititem{What still needs checking}{This step does not yet explain how the
$Y^2$ dependence will be represented in the LFR.}

\subsection{Step 6: Use repeated Delta-block structure to represent richer dependency}

\audititem{Drenth reference}{Thesis Section~2.1, eq.~(2.2), and Section~2.1.1,
eq.~(2.9) \cite{drenth2025thesis}.}

\audititem{What Drenth does}{Drenth defines $\Delta(p)$ as a repeated
block-diagonal structure and uses repeated copies of the same scheduling
variable in the rational MSD example.}

\audititem{Generalized step}{If richer dependency must be represented, reuse the
same scheduling variable repeatedly inside $\Delta(p)$ rather than introducing
more independent scheduling coordinates by default.}

\audititem{Dual-gantry application}{The dual-gantry realization chooses
\begin{equation}
    \Delta(Y) = Y\,I_6.
    \label{eqv:Delta}
\end{equation}
This is repeated-scalar structure: the same physical scheduling variable $Y$
acts repeatedly on the latent loop signals.}

\audititem{Status}{\statusline{Direct from Drenth} for the admissibility of
repeated Delta structure. \statusline{Own derivation} for the specific choice
$YI_6$ and the repetition count $6$.}

\audititem{What still needs checking}{Drenth does not justify why the dimension
should be exactly $6$ or whether this choice is minimal.}

\subsection{Step 7: Introduce latent variables so the Delta action becomes simple}

\audititem{Drenth reference}{Thesis Section~2.1, eq.~(2.1), together with the
rational MSD structure in Section~2.1.1 \cite{drenth2025thesis}.}

\audititem{What Drenth does}{Drenth uses latent variables $z$ and $w$ as the
internal mechanism by which scheduling dependence is routed through
$w = \Delta(p)z$. However, he does not give a symbolic algorithm for choosing
these variables for a new plant.}

\audititem{Generalized step}{Introduce internal algebraic helper variables so
that the scheduling dependence can be represented through simple repeated action
of the Delta block.}

\audititem{Dual-gantry application}{Define the generalized force
\begin{equation}
    f_{\mathrm{gen}} = [-K,\,-C]\,x + u,
    \label{eqv:fgen}
\end{equation}
then define the plant-specific latent variables
\begin{equation}
    v = M(Y)^{-1} f_{\mathrm{gen}}, \qquad
    v_1 = Yv, \qquad
    v_2 = Y^2 v = Y v_1.
    \label{eqv:latent}
\end{equation}
Here $v$, $v_1$, and $v_2$ are whole vectors in $\R^3$. In logical coordinates,
$v$ corresponds to the acceleration vector
\[
v = \ddot q = \begin{bmatrix} \ddot X \\ \ddot \Theta \\ \ddot Y \end{bmatrix},
\]
and the subscripts on $v_1$ and $v_2$ label derived latent vectors rather than
components of $v$. These are internal helper variables, not physical states.}

\audititem{Status}{\statusline{Direct from Drenth} for the latent-variable
concept. \statusline{Own derivation} for the specific choices
$v$, $v_1$, and $v_2$.}

\audititem{What still needs checking}{This is one of the main places where
Drenth stops and our own derivation begins. The actual latent-variable choice is
fully plant-specific.}

\subsection{Step 8: Assemble the loop signals and push all scheduling dependence into the loop}

\audititem{Drenth reference}{Thesis Section~2.1 for the structural target,
reinforced by Section~2.1.1 for repeated scheduling action
\cite{drenth2025thesis}.}

\audititem{What Drenth does}{Drenth's framework requires a constant
interconnection matrix and a Delta loop containing the scheduling dependence.
The MSD example shows repeated scheduling action in such a loop but does not
give a universal symbolic construction rule.}

\audititem{Generalized step}{Package the helper variables into $z$ and $w$ so
that $w=\Delta(p)z$ holds directly and all non-constant dependence is pushed
into the loop.}

\audititem{Dual-gantry application}{Choose
\begin{equation}
    z = \begin{bmatrix} v \\ v_1 \end{bmatrix}, \qquad
    w = \begin{bmatrix} v_1 \\ v_2 \end{bmatrix},
    \label{eqv:zw}
\end{equation}
so that $w = \Delta(Y) z = Y z$. Also decompose the inertia matrix as
\begin{equation}
    M(Y) = M_0 + Y M_1 + Y^2 M_2.
    \label{eqv:decomp}
\end{equation}
Then pre-multiplying $v = M(Y)^{-1} f_{\mathrm{gen}}$ by $M(Y)$ yields
\begin{equation}
    M_0 v = f_{\mathrm{gen}} - M_1 v_1 - M_2 v_2,
    \label{eqv:M0v}
\end{equation}
so the dependence on $Y$ has been shifted into the loop variables.}

\audititem{Status}{\statusline{Generalized from Drenth} for the design
objective. \statusline{Own derivation} for the decomposition of $M(Y)$, the
signal definitions in \eqref{eqv:zw}, and the algebraic rewrite in
\eqref{eqv:M0v}.}

\audititem{What still needs checking}{This is not a direct Drenth construction.
The decomposition of $M(Y)$ and the specific algebraic rewrite are part of the
dual-gantry derivation.}

\subsection{Step 9: Rewrite the plant in constant-matrix LPV-LFR form and read off $G$}

\audititem{Drenth reference}{Thesis Section~2.1, eq.~(2.1)
\cite{drenth2025thesis}.}

\audititem{What Drenth does}{Drenth fixes the block structure that a valid
LPV-LFR must match. The actual coefficients are model-dependent.}

\audititem{Generalized step}{Once the equations are in LPV-LFR form, collect the
coefficients of $x$, $w$, and $u$ and identify the constant blocks of $G$.}

\audititem{Dual-gantry application}{Using the dual-gantry rewrite, the constant
blocks become
\begin{align}
    A &=
    \begin{bmatrix}
        0 & I_3 \\
        -M_0^{-1}K & -M_0^{-1}C
    \end{bmatrix},
    &
    B_w &=
    \begin{bmatrix}
        0 & 0 \\
        -M_0^{-1}M_1 & -M_0^{-1}M_2
    \end{bmatrix},
    &
    B_u &=
    \begin{bmatrix}
        0 \\
        M_0^{-1}
    \end{bmatrix},
    \label{eqv:stateblocks}
\end{align}
\begin{align}
    C_z &=
    \begin{bmatrix}
        M_0^{-1}[-K,\,-C] \\
        0
    \end{bmatrix},
    &
    D_{zw} &=
    \begin{bmatrix}
        -M_0^{-1}M_1 & -M_0^{-1}M_2 \\
        I_3 & 0
    \end{bmatrix},
    &
    D_{zu} &=
    \begin{bmatrix}
        M_0^{-1} \\
        0
    \end{bmatrix},
    \label{eqv:loopblocks}
\end{align}
with
\begin{equation}
    C_y = \begin{bmatrix} I_3 & 0 \end{bmatrix},
    \qquad
    D_{yw} = 0,
    \qquad
    D_{yu} = 0.
    \label{eqv:outblocks}
\end{equation}}

\audititem{Status}{\statusline{Direct from Drenth} for the required block
partition. \statusline{Own derivation} for the explicit matrix formulas.}

\audititem{What still needs checking}{Drenth does not justify uniqueness or
minimality of these formulas. This is plant-specific coefficient matching.}

\subsection{Step 10: Collapse the loop and verify exact recovery of the target CT model}

\audititem{Drenth reference}{Thesis Section~2.1, eqs.~(2.3)--(2.4)
\cite{drenth2025thesis}.}

\audititem{What Drenth does}{Drenth eliminates the latent variables to recover
the collapsed LPV state-space model.}

\audititem{Generalized step}{Use loop collapse as the main exactness check:
substitute the Delta relation, eliminate the latent variables, and verify that
the target continuous-time model is recovered.}

\audititem{Dual-gantry application}{With $w = \Delta(Y) z = Y z$, we have
\[
w_1 = Yv, \qquad w_2 = Yv_1 = Y^2 v.
\]
Substituting into the first block of
$z = C_z x + D_{zw} w + D_{zu} u$ gives
\[
v = M_0^{-1}[-K,-C]\,x - M_0^{-1}M_1\,(Yv) - M_0^{-1}M_2\,(Y^2 v) + M_0^{-1}u.
\]
Multiplying through by $M_0$ and collecting yields
\begin{equation}
    M_0 v + Y M_1 v + Y^2 M_2 v = f_{\mathrm{gen}},
    \label{eqv:loopcollapse}
\end{equation}
hence
\begin{equation}
    M(Y)\,v = f_{\mathrm{gen}}.
    \label{eqv:MYv}
\end{equation}
Then the state equation becomes
\[
\dot{x} = \begin{bmatrix} \dot{q} \\ v \end{bmatrix}
= \begin{bmatrix} \dot{q} \\ M(Y)^{-1} f_{\mathrm{gen}} \end{bmatrix}
= A_c(Y)\,x + B_c(Y)\,u,
\]
which matches \eqref{eqv:ctss}--\eqref{eqv:ctmats}.}

\audititem{Status}{\statusline{Direct from Drenth} for the collapse principle.
\statusline{Own derivation} for the explicit reduction to \eqref{eqv:MYv}.}

\audititem{What still needs checking}{The identity $M(Y)v = f_{\mathrm{gen}}$
is not a Drenth theorem. It is the key plant-specific correctness result of the
dual-gantry derivation.}

\subsection{Step 11: Confirm the collapsed dependence class}

\audititem{Drenth reference}{Thesis Section~2.1, text after eq.~(2.4)
\cite{drenth2025thesis}.}

\audititem{What Drenth does}{Drenth states that collapsing the LPV-LFR yields an
LPV state-space model with rational dependence on the scheduling variable, with
affine dependency as the special case $D_{zw}=0$.}

\audititem{Generalized step}{After collapse, classify the resulting dual-gantry
model as affine or rational.}

\audititem{Dual-gantry application}{Because the collapsed matrices contain
$M(Y)^{-1}$, the dual-gantry CT model is rational in $Y$.}

\audititem{Status}{\statusline{Direct from Drenth} for the classification
framework. \statusline{Own derivation} for classifying this specific plant.}

\audititem{What still needs checking}{This does not say whether an approximate
affine embedding might also be useful. It only classifies the exact baseline.}

%------------------------------------------------------------------------------
\section{Well-Posedness: Generic Drenth Route Versus Sharper Plant-Specific Route}
%------------------------------------------------------------------------------

\subsection{Step 12: State the exact basic loop-solvability condition}

\audititem{Drenth reference}{Thesis Section~2.1, immediately after eq.~(2.4)
\cite{drenth2025thesis}.}

\audititem{What Drenth does}{Drenth states the exact algebraic-loop solvability
condition:
\[
I - D_{zw}\Delta(p)
\]
must be nonsingular for all admissible scheduling values.}

\audititem{Generalized step}{Every candidate LPV-LFR realization must first be
checked against the exact basic nonsingularity condition.}

\audititem{Dual-gantry application}{For the dual-gantry realization, the exact
condition is
\begin{equation}
    I - D_{zw}\Delta(Y) \text{ is nonsingular for all admissible } Y.
    \label{eqv:wpbasic}
\end{equation}}

\audititem{Status}{\statusline{Direct from Drenth}.}

\audititem{What still needs checking}{At this stage we do not yet have either
Drenth's sufficient theorem or the sharper plant-specific simplification. Those
are two different next moves and must be kept separate.}

\subsection{Step 13: State Drenth's generic sufficient well-posedness route}

\audititem{Drenth reference}{Thesis Section~2.2, Assumptions~2.1--2.2,
Remark~2.3, Condition~2.4, Theorem~2.5, and eqs.~(2.10)--(2.13)
\cite{drenth2025thesis}.}

\audititem{What Drenth does}{Drenth turns the exact basic condition into a
generic sufficient route:
\begin{itemize}[leftmargin=1.2em]
    \item $\Delta(p)$ has diagonal structure,
    \item the scheduling set is contained in the unit $\ell_\infty$ ball,
    \item $\rho(D_{zw}) < 1$,
    \item therefore $I - D_{zw}\Delta(p)$ is nonsingular for all admissible $p$.
\end{itemize}
He then introduces parameterizations of $D_{zw}$ suited for unconstrained
optimization.}

\audititem{Generalized step}{One generic way to establish well-posedness is:
\begin{enumerate}[leftmargin=1.4em,label=\arabic*.]
    \item verify the required Delta structure,
    \item normalize the scheduling range if needed,
    \item verify the spectral-radius bound on $D_{zw}$,
    \item apply Theorem~2.5.
\end{enumerate}
This route is sufficient, not exact.}

\audititem{Dual-gantry application}{For the dual gantry:
\begin{itemize}[leftmargin=1.2em]
    \item the chosen $\Delta(Y)=YI_6$ is structurally compatible with the
    diagonal repeated-block requirement,
    \item on the operational interval $Y \in [0.05,0.75]$, the raw variable
    already satisfies $|Y| \leq 1$,
    \item for another bounded interval, Remark~2.3 says scaling and shifting are
    available,
    \item the unresolved part is Condition~2.4: we would still need to prove
    $\rho(D_{zw}) < 1$ for the specific dual-gantry $D_{zw}$.
\end{itemize}
So Drenth's Section~2.2 gives a generic sufficient route in principle, but that
route is not completed for the present gantry derivation unless the spectral
radius is checked explicitly. In other words, this route is available as theory
background, but it is not yet the proof route actually used for the current
baseline.}

\audititem{Status}{\statusline{Direct from Drenth} for the theorem and its
assumptions. \statusline{Generalized from Drenth} for the workflow
``check assumptions, then apply theorem.'' \statusline{Own derivation} for
identifying how the current dual-gantry realization fits those assumptions.}

\audititem{What still needs checking}{This route does not by itself justify
global well-posedness for all $Y \in \R$. It also does not justify invoking
Theorem~2.5 without actually checking $\rho(D_{zw}) < 1$.}

\subsection{Step 14: State the sharper plant-specific well-posedness route}

\audititem{Drenth reference}{The exact starting point comes from Section~2.1:
well-posedness is equivalent to solvability of the algebraic loop. The actual
reduction below is \textbf{not} from Drenth; it comes from the dual-gantry
algebra in \texttt{LPV/LFR-derivation.tex} and the companion proof in
\texttt{LPV/M-invertibility.tex}.}

\audititem{What Drenth does}{Drenth supplies the exact loop-solvability
condition and a generic sufficient theorem, but not a plant-specific reduction
for this mechanical model.}

\audititem{Generalized step}{If the plant structure allows it, reduce the exact
loop-solvability condition to a simpler condition that is specific to the plant
and sharper than the generic sufficient theorem.}

\audititem{Dual-gantry application}{For the dual gantry, the loop equations are
\begin{align}
    v   &= M_0^{-1}[-K,-C]\,x - M_0^{-1}M_1\,(Yv)
           - M_0^{-1}M_2\,(Yv_1) + M_0^{-1}u,
           \label{eqv:loopv} \\
    v_1 &= Yv. \label{eqv:loopv1}
\end{align}
Equivalently, starting from Drenth's generic loop relation
\[
    (I-D_{zw}\Delta(Y))\,z = C_z x + D_{zu}u,
\]
and using
\[
    z=\begin{bmatrix}v\\v_1\end{bmatrix},\qquad
    D_{zw}=
    \begin{bmatrix}
        -M_0^{-1}M_1 & -M_0^{-1}M_2 \\
        I_3 & 0
    \end{bmatrix},
    \qquad
    \Delta(Y)=YI_6=
    \begin{bmatrix}
        YI_3 & 0 \\
        0 & YI_3
    \end{bmatrix},
\]
gives
\[
    I-D_{zw}\Delta(Y)=
    \begin{bmatrix}
        I_3 + YM_0^{-1}M_1 & YM_0^{-1}M_2 \\
        -YI_3 & I_3
    \end{bmatrix},
\]
so the exact loop equation becomes
\[
    \begin{bmatrix}
        I_3 + YM_0^{-1}M_1 & YM_0^{-1}M_2 \\
        -YI_3 & I_3
    \end{bmatrix}
    \begin{bmatrix}
        v \\
        v_1
    \end{bmatrix}
    =
    \begin{bmatrix}
        M_0^{-1}[-K,-C]\,x + M_0^{-1}u \\
        0
    \end{bmatrix}.
\]
This is equivalent to the two block equations \eqref{eqv:loopv}--\eqref{eqv:loopv1}. The
second block gives \(v_1=Yv\), and substituting \eqref{eqv:loopv1} into
\eqref{eqv:loopv} yields
\[
(M_0 + YM_1 + Y^2 M_2)\,v = f_{\mathrm{gen}},
\]
that is,
\begin{equation}
    M(Y)\,v = f_{\mathrm{gen}}.
    \label{eqv:loopreduced}
\end{equation}
So the algebraic loop has a unique solution if and only if $M(Y)$ is
invertible.

The companion note \texttt{LPV/M-invertibility.tex} then proves, using
Sylvester's criterion, that under positive masses, inertias, and geometric
parameters:
\[
M(Y) \succ 0 \qquad \forall\, Y \in \R.
\]
Hence $M(Y)$ is invertible for all real $Y$, and therefore the dual-gantry LFR
is globally well-posed in $Y$. This is the proof route actually used for the
current baseline.}

\audititem{Status}{\statusline{Direct from Drenth} only for the fact that the
exact algebraic loop must be solved. \statusline{Own derivation} for the
reduction to \eqref{eqv:loopreduced}. \statusline{Own derivation} for the proof
route through the companion invertibility note \texttt{LPV/M-invertibility.tex}.}

\audititem{What still needs checking}{This route should not be described as an
application of Theorem~2.5. It is a different, sharper route.
Also, if physical parameters are later made trainable, the proof remains valid
only if the positivity assumptions in the companion invertibility note are
preserved.}

\subsection{Step 15: Compare the two well-posedness routes explicitly}

\audititem{Drenth reference}{Section~2.1 gives the exact condition, and
Section~2.2 gives a sufficient theorem \cite{drenth2025thesis}.}

\audititem{What Drenth does}{Drenth does not perform this comparison for the
dual gantry. The comparison here is a verification aid built on top of the
framework he provides.}

\audititem{Generalized step}{When both a generic sufficient theorem and a
plant-specific exact reduction are available, compare them rather than treating
them as interchangeable.}

\audititem{Dual-gantry application}{For the dual gantry, the two routes are best
read side by side:
\begin{center}
\begin{tabular}{p{0.22\linewidth}p{0.33\linewidth}p{0.33\linewidth}}
\toprule
Aspect & Drenth Section~2.2 route & Dual-gantry plant-specific route \\
\midrule
Type of result & Generic sufficient theorem & Exact structural reduction for this plant \\
Starting point & Diagonal $\Delta$, bounded scheduling, $\rho(D_{zw})<1$ & Exact loop equations of the chosen realization \\
Key object to check & Spectral radius of $D_{zw}$ & Invertibility of $M(Y)$ \\
Range in which it naturally lives & Bounded scheduling setup & Global in $Y$ once $M(Y)\succ 0$ is proved \\
Best use here & Background theorem and later learnable-model context & Actual baseline well-posedness proof \\
\bottomrule
\end{tabular}
\end{center}
So the honest conclusion is:
\begin{itemize}[leftmargin=1.2em]
    \item Drenth Section~2.2 supplies the \textbf{generic sufficient theorem}.
    \item The dual-gantry derivation supplies the \textbf{sharper plant-specific
    proof}.
\end{itemize}}

\audititem{Status}{\statusline{Generalized from Drenth} for the comparison
viewpoint. \statusline{Own derivation} for concluding which route is sharper for
this plant.}

\audititem{What still needs checking}{This comparison does not make Drenth's
route irrelevant. It remains useful if one later introduces a learnable rational
LPV-LFR block or wants a bounded-range sufficient theorem aligned with Drenth's
identification framework.}

%------------------------------------------------------------------------------
\section{Boundary Summary}
%------------------------------------------------------------------------------

\begin{resultbox}{Directly supported by Drenth}
\begin{itemize}[leftmargin=1.2em]
    \item CT LPV-LFR framework $(G,\Delta(p))$
    \item repeated block-diagonal Delta structure
    \item latent-loop interpretation of $z$ and $w$
    \item affine special case $D_{zw}=0$
    \item collapse of the loop to an LPV-SS model
    \item rational dependency after collapse
    \item affine-versus-rational trade-off and overbounding motivation
    \item exact basic loop-solvability condition
    \item generic sufficient well-posedness theorem from Section~2.2
    \item optimization-oriented well-posed parameterization of $D_{zw}$
\end{itemize}
\end{resultbox}

\begin{resultbox}{Generalized from Drenth}
\begin{itemize}[leftmargin=1.2em]
    \item using the LPV-LFR form as a design target for a plant-specific reformulation
    \item using collapse as the main correctness check of a candidate realization
    \item using coupling preservation as a scheduling-design principle
    \item preferring repeated use of one physical scheduling variable over
    introducing unnecessary independent scheduling variables
    \item treating Theorem~2.5 as one available well-posedness route rather than
    the only route
\end{itemize}
\end{resultbox}

\begin{resultbox}{Own Dual-Gantry Derivation}
\begin{itemize}[leftmargin=1.2em]
    \item the target CT gantry model itself
    \item the decomposition $M(Y) = M_0 + YM_1 + Y^2M_2$
    \item the specific latent variables $v$, $v_1$, $v_2$
    \item the signal definitions $z = [v;v_1]$, $w = [v_1;v_2]$
    \item the specific Delta block $\Delta(Y) = YI_6$
    \item the explicit constant matrix blocks of $G$
    \item the reduction of the loop to $M(Y)v = f_{\mathrm{gen}}$
    \item the sharper plant-specific well-posedness proof through invertibility
    of $M(Y)$
\end{itemize}
\end{resultbox}

%------------------------------------------------------------------------------
\section{Manual Checking Priorities}
%------------------------------------------------------------------------------

\begin{warningbox}{Weakest or Most Important Parts to Check Manually}
The current document is structured and explicit, but the following parts still
deserve the most manual attention:
\begin{enumerate}[leftmargin=1.4em,label=\arabic*.]
    \item \textbf{Equation numbering against the thesis.} The document assumes
    the current local thesis edition matches the Section~2 numbering already
    extracted in the source-only notes. This should be checked once more against
    the PDF before any external use.
    \item \textbf{The plant-specific latent-variable choice.} The definitions of
    $v$, $v_1$, $v_2$, $z$, and $w$ are the most important own-derivation step.
    They are algebraically sensible, but they are not prescribed by Drenth.
    \item \textbf{The reduction to $M(Y)v = f_{\mathrm{gen}}$.} This is the key
    exactness result of the gantry derivation. It should be checked line by line
    against \texttt{LPV/LFR-derivation.tex}.
    \item \textbf{The relationship between Drenth's Theorem~2.5 and the current
    gantry proof.} The present document keeps them separate, which is good, but
    if the supervisor-facing version later cites Theorem~2.5 directly, it should
    be made explicit that the sharper proof used here is different.
    \item \textbf{Future trainable-parameter use.} If physical parameters are
    later updated during augmentation, the global invertibility statement remains
    valid only if the positivity assumptions from \texttt{LPV/M-invertibility.tex}
    are preserved by the parameterization.
\end{enumerate}
\end{warningbox}

\begin{resultbox}{Practical Outcome}
This verification document now supports the following clean statement:

Drenth's thesis provides the generic CT LPV-LFR framework, the affine-versus-rational
motivation, and a generic sufficient well-posedness theorem. The exact
dual-gantry realization, however, is mainly our own derivation, and the
strongest well-posedness proof for this baseline is the plant-specific reduction
to $M(Y)v = f_{\mathrm{gen}}$ together with the companion invertibility proof.
\end{resultbox}

\newpage
\begin{thebibliography}{9}

\bibitem{garcia2013model}
I.~Garcia-Herreros, X.~Kestelyn, J.~Gomand, R.~Coleman, and P.-J.~Barre,
``Model-based decoupling control method for dual-drive gantry stages:
A case study with experimental validations,''
\textit{Control Engineering Practice}, vol.~21, no.~3, pp.~298--307, 2013.

\bibitem{drenth2025lpvlfr}
J.~Drenth, J.~Hoekstra, M.~Schoukens, and R.~Toth,
``Efficient Learning of Affine and Rational Dependency LPV Models With LFR,''
in \textit{Proc.\ IFAC}, 2025.

\bibitem{drenth2025thesis}
J.~Drenth,
``Gradient-based Learning of LPV-LFR Models,''
M.Sc.\ thesis, Eindhoven University of Technology, 2025.

\end{thebibliography}

\end{document}
